% ===================================================================
%  કોષ્ટક નં. ૭ — શિયાળામાં ગામ. --- વસંતમાં ખેતરનું ઘર.
%
%  Traduction, faite en 2026, en gujarati d'aujourd'hui, sur l'ido de
%  texto/io. Regles de la colonne : voir 10-kostak-01.tex. Deux
%  scenes, deux numerotations. Ouverture de la DEUXIEME SERIE :
%  « બીજી શ્રેણી ». Un bloc marque « ¶2 » par ordo.py, t07-c1-07-1 :
%  une seule cle, deux alineas.
%
%  LES BETES DE LA FERME SONT SANSKRITES, MAIS LA VACHE EST AUTRE
%  CHOSE. ઘોડો le cheval, ઘોડી la jument, બળદ le boeuf, ઘેટું le
%  mouton, બકરી la chevre, કૂતરો le chien, મરઘી la poule, કબૂતર le
%  pigeon, સસલું le lapin : tous d'un fonds indo-aryen que le hindi
%  et le marathi partagent mot pour mot. Mais « ગાય » la vache (25,
%  56) n'est pas un animal de ferme comme les autres dans cette
%  langue-la, et le releve doit le dire une fois : c'est le seul
%  animal du livret qui soit sacre chez ses lecteurs. La colonne ne
%  change rien — la planche montre une vache qu'on trait devant
%  l'etable, et c'est ce qu'on ecrit. Mais « ગોધન », la richesse en
%  vaches, et « ગોશાળા », l'etable, sont des mots que le gujarati
%  emploie encore, et « ગોકુળ » est un nom de lieu. Une langue peut
%  nommer une bete ordinairement et l'entourer de mots qui ne le sont
%  pas.
%
%  ET LE COCHON (77, 78) EST LA VRAIE DIFFICULTE DE CE TABLEAU, comme
%  la charcuterie l'etait au tableau 15 haoussa. Le gujarati est
%  massivement vegetarien — le Gujarat est le seul Etat de l'Inde ou
%  la majorite le soit —, et le porc y est ce qu'il est aussi pour un
%  musulman. Le mot existe pourtant et n'a rien d'obscur : « ભૂંડ »,
%  et « ભૂંડણ » pour la truie. La colonne l'ecrit sans note et sans
%  detour, exactement comme la haoussa avait ecrit naman alade. LA
%  LANGUE A LES MOTS ; C'EST LE MANGER QUI EST INTERDIT, PAS LE
%  NOMMER — et c'est maintenant la deuxieme colonne du releve a le
%  demontrer, dans une autre religion.
%
%  LE LABOUR (5, 8) EST LE PLUS VIEUX MOT DE CETTE COLONNE. « હળ » la
%  charrue vient du sanskrit हल, il est dans les Vedas, et l'objet
%  grave sur la planche — un soc, un age, deux mancherons, un cheval
%  attele — est celui que le mot designe depuis trois mille ans. La
%  colonne persane avait trouve son mot le plus vieux dans le corps,
%  la haoussa dans la basse-cour ; celle-ci le trouve dans le champ.
%
%  LE ROUET DE LA FILEUSE (35) EST « રેંટિયો », ET IL FAUT LE DIRE.
%  C'est le rouet a filer le coton, et c'est l'objet que Gandhi a mis
%  au centre du drapeau de l'Inde avant qu'on n'y mette la roue
%  d'Ashoka ; le charkha de la desobeissance, file une heure par
%  jour. La planche grave une vieille femme qui file le lin ou le
%  chanvre au coin d'une ferme europeenne, en 1926, l'annee meme ou
%  le charkha etait un mot d'ordre. Le mot est ecrit sans hesiter,
%  parce que c'est le meme objet — mais aucune autre colonne du
%  releve ne porte un outil qui ait ete, la meme annee, un drapeau.
%
%  ET LE MARRON DU TABLEAU 7 HAOUSSA A ICI SON PENDANT EXACT : LA
%  CHATAIGNE DE L'ALINEA 18 (42). Le haoussa avait refuse « gyaɗa »,
%  l'arachide grillee, dont la scene est mot pour mot celle de la
%  planche. Le gujarati a « શેકેલી સિંગ », l'arachide grillee qu'on
%  vend en cornet de papier a chaque coin de rue de Ahmedabad l'hiver
%  — meme scene, meme papier, meme saison. On ecrit « ચેસ્ટનટ ».
%  DEUX COLONNES, DEUX ARACHIDES, DEUX REFUS AU MEME NUMERO.
% ===================================================================

%%K t07-noto-1 noto
\VUnotes{143.2mm}{%
(*) ભોજનની વિગત માટે કોષ્ટક ૬ અને ૧૩ જુઓ.%
}

%%K t07-apar-1 apar
\VUsaut{4.0mm}
\VUcentre{13.2pt}{90}{બીજી શ્રેણી}
\VUsaut{3.0mm}
\VUfilet{16mm}
\VUsaut{5.6mm}
\VUtitre{15.0pt}{60}{કોષ્ટક નં. ૭}
\VUsaut{3.0mm}

%%K t07-tit-1 sub
\VUcentre{12.0pt}{0}{\textit{પ્રથમ દૃશ્ય.}}
\VUsaut{3.0mm}

%%K t07-tit-2 sub
\VUpk{10.2pt}{40}{શિયાળામાં ગામ.}
\VUsaut{4.0mm}

%%K t07-c1-01-1 p
1. --- નવા વર્ષની રજાઓમાં, હું મારા પિતાની સાથે નોવુર્બો ગયો, એક નાનું
ગામ જ્યાં તેમને કેટલાંક વેપારી કામ પતાવવાનાં હતાં.

%%K t07-c1-02-1 p
2. --- અમે \VUgras{ઘોડાગાડી} \textsuperscript{(11)} વાપરી, જે શહેર અને
આ ગામ વચ્ચે ચાલે છે ; પણ ઘણી ઠંડી હોવાથી, થોડી અગવડવાળા વાહનમાંની આ
મુસાફરી બહુ સુખદ ન હતી.

%%K t07-c1-03-1 p
3. --- તેથી અમને સાચો આનંદ થયો જ્યારે અમે \VUgras{ગાડીવાનને} પોતાની
ગાડી ચોકમાં ઊભી રાખતાં જોઈ, \VUgras{ધર્મશાળા} \textsuperscript{(2)}
આગળ, જે \textit{સફેદ રીંછ}ની છે. \VUgras{ધર્મશાળાવાળો}
\textsuperscript{(4)} અમારી રાહ જોતો પોતાના ઘરના ઉંબરે ઊભો હતો, શાંતિથી
પોતાની \VUgras{ચલમ} પીતો.

%%K t07-c1-03-2 p
તેણે અમને અંદર આવવા કહ્યું અને હું વિનંતીની રાહ જોયા વિના સગડીમાં
ભડભડતી ડાળીઓની આગ આગળ બેસી ગયો. ત્યારે મેં તીખા \VUgras{ઉત્તરના
પવનની} ખૂબ મજાક કરી, જે જૂના \VUgras{પાટિયાને} \textsuperscript{(3)}
કિચૂડાટ કરાવે છે અને જે સગડીમાંથી નીકળતા \VUgras{ધુમાડાને}
\textsuperscript{(1)} કાળા વમળોમાં ઉડાડી લઈ જાય છે.

%%K t07-c1-04-1 p
4. --- બપોરના ભોજનનો સમય હતો. વધારે રાહ જોયા વિના, અમે અમને પીરસાયેલું
સાદું પણ સ્વાદિષ્ટ ભોજન સારી રીતે જમ્યા (*).

%%K t07-tit-3 sub
\VUpk{10.2pt}{40}{કેટલીક મુલાકાતો.}
\VUsaut{3.0mm}

%%K t07-c1-05-1 p
5. --- બપોરના ભોજન પછી, હું મારા પિતાની સાથે ગયો, જે વીમા એજન્ટ છે,
તેમના ગ્રાહકોને ત્યાં. પહેલાં અમે \VUgras{લુહારને} \textsuperscript{(5)}
મળ્યા, જે ધર્મશાળા પાસે રહે છે. તે ઘોડાને નાળ જડવામાં રોકાયેલો હતો.
મેં તેના સાથીની તાકાતની પ્રશંસા કરી, જેણે ઘોડાની \VUgras{ખરી} પોતાના
ચામડાના એપ્રન પર મજબૂત રીતે પકડી રાખી, જ્યારે લુહાર જોરદાર હથોડા મારીને
\VUgras{નાળ} \textsuperscript{(6)} જડતો હતો.

%%K t07-c1-06-1 p
6. --- પછી અમે ગામના \VUgras{મોચી} \textsuperscript{(7)} પાસે ગયા,
વૃદ્ધ યાકોબુસ. જોકે તેની પાસે નિશાની તરીકે ઘણો સુંદર \VUgras{બૂટ} છે,
તે કાણાં પડેલાં જૂનાં ચંપલને નવાં તળિયાં લગાડવાથી સંતોષ માને છે.

%%K t07-c1-06-2 p
વૃદ્ધે અમને હસતાં હસતાં કહ્યું : « શહેરમાં હું « ચંપલ બનાવનારો » હતો
અને મારે ગ્રાહકો ન હતા. અહીં હું « ચંપલ સાંધનારો » છું અને કામ ઘણું છે. »

%%K t07-c1-07-1 p
7. --- તેણે અમને પોતાના પડોશી \VUgras{હજામ} \textsuperscript{(8)} વિશે
વાત કરી, જેના ગ્રાહકો સંતુષ્ટ નથી. તેના \VUgras{અસ્તરા} હંમેશાં બુઠ્ઠા
હોય છે અને તેની \VUgras{કાતર} કદી બરાબર કાપતી નથી.

પણ તે ગામનો એકમાત્ર હજામ હોવાથી, સ્વાભાવિક રીતે તેની પાસે જ હજામત
કરાવવી અને વાળ કપાવવા પડે છે. વળી તે કંજૂસ અને અણગમતો માણસ છે.

%%K t07-c1-08-1 p
8. --- સદ્ભાગ્યે, બીજા \VUgras{પડોશીઓ} તેના જેવા નથી. દાખલા તરીકે
\VUgras{ગાડી બનાવનારો} \textsuperscript{(9)} ભલો માણસ છે, મહેનતુ અને
મદદરૂપ. તેની પાસે ગાડી સમરાવવી એ આનંદ છે. તેનું કામ હંમેશાં પૂરું
થયેલું હોય છે.

%%K t07-c1-09-1 p
9. --- વૃદ્ધ યાકોબુસે \VUgras{પલાણ બનાવનારા} \textsuperscript{(10)}
વિશે પણ ફરિયાદ કરી નહીં, જેને તે ક્યારેક કામ ઉતાવળનું હોય ત્યારે
\VUgras{ઘોડાનો સામાન} સીવવામાં મદદ કરે છે.

%%K t07-c1-09-2 p
મારા પિતાએ તેને તેના કુટુંબ વિશે પૂછ્યું : « બધાં સારાં છે. મારી
પૌત્રી \VUgras{શાળામાં} \textsuperscript{(12)} છે. તેની
\VUgras{શિક્ષિકા} \textsuperscript{(13)} તેનાથી ઘણી સંતુષ્ટ છે, તે સારી
\VUgras{વિદ્યાર્થિની} \textsuperscript{(14)} છે. તે મારા ખરાબ છોકરા
જેવી નથી ; તે માત્ર રમવાનું જ વિચારે છે. \VUgras{શિક્ષક}
\textsuperscript{(15)} તેને કશું શીખવી શકતા નથી. »

%%K t07-tit-4 sub
\VUsaut{3.0mm}
\VUpk{10.2pt}{40}{ગામની વાતો.}
\VUsaut{3.0mm}

%%K t07-c1-10-1 p
10. --- પછી વૃદ્ધે અમને ગામના સમાચાર કહ્યા. નવી શાળા બંધાવાની છે, કેમ
કે \VUgras{પંચાયતના ઘરમાં} બેસાડેલી શાળા બહુ નાની પડી છે. વળી એ સહેલું
છે, કેમ કે \VUgras{ઘંટીવાળાના} બગીચાનો એક ભાગ ખરીદવો બસ થશે. એ બિચારા
માણસને ઘણો ફાયદો થશે. ઠંડીને લીધે, \VUgras{ઘંટીના} \textsuperscript{(16)}
\VUgras{પડ} ઘઉં દળી શકતાં નથી, કેમ કે \VUgras{ચક્કર}
\textsuperscript{(17)} થંભી ગયું છે, જેને \VUgras{બરફે}
\textsuperscript{(25)} થિજાવી દીધું છે, જે \VUgras{નાની નદીનો}
\textsuperscript{(23)} છે.

%%K t07-c1-11-1 p
11. --- « બાંધકામની વાત કરીએ તો, તમે અમારું નવું \VUgras{દેવળ}
\textsuperscript{(18)} જોયું ? તે પોતાની બરફની ચાદર નીચે ઘણું મનોહર છે.
\VUgras{કાગડા} \textsuperscript{(19)}, જે પોતાનું જૂનું ઘંટમંદિર હવે
ઓળખતા નથી, \VUgras{કબ્રસ્તાનના} \textsuperscript{(20)} મોટા ઝાડ પર
ઉદાસ બેસે છે. »

%%K t07-c1-12-1 p
12. --- કબ્રસ્તાનના આ વિચારે મોચીને એ માણસો યાદ કરાવ્યા જેમને મારા
પિતા ઓળખતા હતા અને જે ગયા વર્ષે મરી ગયા. « કેટલી \VUgras{કબરો}
\textsuperscript{(21)}, મારા સાહેબ, બીજી કબરોમાં ઉમેરાઈ,
\VUgras{સરુના ઝાડ} \textsuperscript{(22)} નીચે ! મોતે અમારા જૂના
નોટરીને લણી લીધા. આખું ગામ તેમની \VUgras{સ્મશાનયાત્રામાં} હાજર રહ્યું. »

%%K t07-c1-13-1 p
13. --- « આ રહ્યા તેમના અનુગામી, જે નાની નદી પર લસરે છે. તે હમણાં જ
મારી પાસે પોતાની \VUgras{લસરવાની જોડ} \textsuperscript{(26)} નો પટ્ટો
સમરાવવા આવ્યા હતા, જે તૂટી ગયો હતો. »

%%K t07-c1-14-1 p
14. --- « આ રહ્યો નાની નદીના કાંઠે પણ, નવો \VUgras{ગામનો ચોકીદાર}
\textsuperscript{(27)}. તે ભૂતપૂર્વ સિપાહી છે, જે ગામના તોફાની
છોકરાઓને ડરાવે છે. તેની \VUgras{પગની ચામડાની પટ્ટીઓ}
\textsuperscript{(29)} અને તેની \VUgras{ટોપી} \textsuperscript{(28)}
દેખાતાં જ તેઓ ભાગી જાય છે. »

%%K t07-c1-15-1 p
15. --- « અરે ! આ રહ્યો વૃદ્ધ યોસેફ જે ભઠિયારા માટે \VUgras{બળતણના
ભારાની નાની ગાડી} \textsuperscript{(30)} હાંકે છે. --- એ સાચું, મારા
પિતાએ કહ્યું, હું તેની \VUgras{ખચ્ચરની} \textsuperscript{(31)} જોડ
ઓળખું છું. મારે તેની સાથે વાત કરવી છે, હું આ તક લઈશ. ફરી મળીશું, બાપા
યાકોબુસ. »

%%K t07-c1-16-1 p
16. --- અમે નીકળ્યા તે જ વખતે, ગામનાં છોકરાં નિશાળ છોડીને દોડતાં
\VUgras{લસરવાની જગ્યાએ} \textsuperscript{(33)} ધસી ગયાં, પડવાનું અને
\VUgras{પડવામાં} \textsuperscript{(34)} અંગ ભાંગવાનું જોખમ વહોરીને.
તેમાંના એકે ગાડી બનાવનારા પાસેથી પોતાની \VUgras{લસરગાડી}
\textsuperscript{(32)} લીધી, તેમાં પોતાના એક સાથીને બેસાડ્યો, અને દોડતો
ઊપડ્યો.

%%K t07-c1-17-1 p
17. --- બીજાં \VUgras{બરફના ઢગલા} \textsuperscript{(37)} તરફ ધસી ગયાં,
જે \VUgras{પુલ} \textsuperscript{(40)} પાસે છે, અને તેમણે
\VUgras{બરફના પૂતળા} \textsuperscript{(39)} પર મારો ચલાવ્યો, જે તેમણે
ગઈ કાલે બનાવ્યું હતું અને જેને તેમણે ટોપી, ચલમ અને ખભે ભરાવેલી નિશાળની
થેલી પહેરાવી હતી.

%%K t07-c1-18-1 p
18. --- તેમને થિજાવતા પવનની અને ઠંડીની બહુ પરવા નથી. મોટા ભાગનાં પાસે
\VUgras{ગળાનો પટ્ટો} \textsuperscript{(35)} પણ નથી અને,
\VUgras{બરફના ગોળાની} \textsuperscript{(38)} લડાઈ પછી ફરી ગરમ થવા માટે,
તેમને મુઠ્ઠીભર ઘણી ગરમ \VUgras{ચેસ્ટનટ} \textsuperscript{(42)} બસ થાય
છે, જે વૃદ્ધ \VUgras{વેચનારી} યાકોબિના પાસેથી ખરીદેલી હોય છે, જે
પોતાની બહુ પાતળી \VUgras{ઓઢણી} \textsuperscript{(43)} નીચે ઠંડીથી
ધ્રૂજે છે.

%%K t07-tit-5 sub
\VUsaut{3.0mm}
\VUcentre{12.0pt}{0}{\textit{બીજું દૃશ્ય.}}
\VUsaut{3.0mm}

%%K t07-tit-6 sub
\VUpk{10.2pt}{40}{વસંતમાં ખેતરનું ઘર.}
\VUsaut{4.0mm}

%%K t07-c2-01-1 p
1. --- શિયાળાના બધા આનંદ છતાં, અમે \VUgras{વસંતના} પાછા ફરવાને ઊંડા
હર્ષથી વધાવીએ છીએ.

%%K t07-c2-01-2 p
મે મહિનામાં, સાંજે, ખેતરનું આંગણું કેવું આનંદદાયક ચિત્ર લાગે છે !

%%K t07-c2-02-1 p
2. --- ક્ષિતિજે \VUgras{સૂરજ} \textsuperscript{(1)} ધીમે ધીમે આથમે છે,
\VUgras{ગામડાને} \textsuperscript{(3)} પોતાનાં જાંબુડિયાં
\VUgras{કિરણોથી} \textsuperscript{(2)} ભરી દેતો. મહેનતુ \VUgras{ખેડૂત}
\textsuperscript{(6)} પોતાનું \VUgras{ખેતર} \textsuperscript{(4)} ખેડવાની
ઉતાવળ કરે છે.

%%K t07-c2-02-2 p
પોતાના \VUgras{હળથી} \textsuperscript{(5)}, જે જોરદાર \VUgras{ઘોડા}
\textsuperscript{(7)} સાથે જોડ્યું છે, તે \VUgras{ચાસ}
\textsuperscript{(8)} પાડે છે જેમાં \VUgras{વાવનારો}
\textsuperscript{(9)} ભરી મુઠ્ઠીથી \VUgras{બી} નાખશે, જે સોનેરી ડૂંડાં
આપશે.

%%K t07-c2-03-1 p
3. --- \VUgras{કાપનારાઓએ} \textsuperscript{(11)} પોતાનાં દાતરડાંથી
ઘાસના મેદાનનું ઘાસ કાપ્યું, ઘાસ સૂકવનારાઓએ \VUgras{સૂકું ઘાસ}
\textsuperscript{(10)} \VUgras{પંજેટીથી} \textsuperscript{(12)} અને
દાંતાળીથી ફેરવીને સૂકવ્યું, અને આ રહ્યા તેઓ પાછા ફરે છે.

%%K t07-c2-03-2 p
કેટલાક બળદ ખેંચતા ભારે ગાડા આગળ ચાલે છે ; તેઓ પોતાનું ઓજાર ખભે લઈ
જાય છે. બીજા, વધારે આળસુ, સુગંધી સૂકા ઘાસ પર આરામથી ગોઠવાઈ ગયા છે.

%%K t07-c2-04-1 p
4. --- પસાર થતાં તેઓ \VUgras{માળીને} \textsuperscript{(16)} મિત્રભાવે
નમસ્કાર કરે છે, જે \VUgras{ફળના બગીચામાં} \textsuperscript{(13)}
\VUgras{ફળઝાડ} છાંટવામાં અને \VUgras{નાસપતીનાં ઝાડ}
\textsuperscript{(14)} ને કલમ કરવામાં રોકાયેલો છે. \VUgras{આલૂબખારાનાં
ઝાડ} \textsuperscript{(15)} ફૂલવા લાગ્યાં છે, અને, સારી રીતે ખાતર
પાયેલા \VUgras{શાકના બગીચામાં} \textsuperscript{(17)}, શાકભાજી, કોબી અને
સલાડ પહેલેથી સારાં થયાં છે.

%%K t07-c2-04-2 p
નાની ભીંત પર, સુંદર \VUgras{મોર} \textsuperscript{(18)} પોતાની પૂંછડી
ફેલાવે છે, જેથી તેનાં ઘણાં સુંદર પીંછાંની પ્રશંસા થાય.

%%K t07-tit-7 sub
\VUpk{10.2pt}{40}{ખેતરનું આંગણું.}
\VUsaut{3.0mm}

%%K t07-c2-05-1 p
5. --- \VUgras{ઘોડારનાં} \textsuperscript{(19)} બારણાં ખુલ્લાં છે અને
ઘાસદાની તથા \VUgras{પથારી} \textsuperscript{(23)} દેખાય છે, જે
\VUgras{ઘોડીની} \textsuperscript{(21)} છે, જેને \VUgras{સાઈસે}
\textsuperscript{(20)} હમણાં જ છોડી. \VUgras{વછેરું}
\textsuperscript{(22)}, પોતાના લાંબા પગ સાથે એટલું રમૂજી, કૂદતું કૂદતું
પોતાની માની પાછળ જાય છે.

%%K t07-c2-06-1 p
6. --- \VUgras{કૂવા} \textsuperscript{(29)} પાસે \VUgras{ગાયો}
\textsuperscript{(25)} લાકડાની \VUgras{કૂંડીમાં} \textsuperscript{(26)}
પાણી પીવા જાય છે, જે તેમના માટે અવેડો છે. \VUgras{વાછરડું}
\textsuperscript{(24)} આ તકનો લાભ લઈને ધાવે છે, અને \VUgras{બળદ}
\textsuperscript{(27)} ઘાસની સારી સુગંધ સૂંઘતો આનંદથી ભાંભરે છે અને
પોતાનું માથું પ્રબળ \VUgras{શિંગડાં} \textsuperscript{(28)} સાથે
ગર્વથી ઊંચકે છે.

%%K t07-c2-07-1 p
7. --- બધાં કામ કરે છે. \VUgras{નોકરાણી} \textsuperscript{(32)} હાથમાં
કૂવાની \VUgras{દોરી} \textsuperscript{(31)} પકડે છે. તે \VUgras{ડોલથી}
\textsuperscript{(30)} પાણી ખેંચે છે ઢોરને પાવા માટે. \VUgras{ખળાની}
\textsuperscript{(33)} પાસે એક માણસ પરાળ ભેગું કરે છે અને \VUgras{ઘરના}
\textsuperscript{(34)} બારણા આગળ વૃદ્ધ \VUgras{કાંતનારી}
\textsuperscript{(35)} પોતાના \VUgras{રેંટિયા} અને પોતાની
\VUgras{ત્રાકથી} \VUgras{શણ} કે \VUgras{ગાંજો} કાંતે છે, જૂના જમાનાની
જેમ.

%%K t07-tit-8 sub
\VUsaut{3.0mm}
\VUpk{10.2pt}{40}{ખેડૂત.}
\VUsaut{3.0mm}

%%K t07-c2-08-1 p
8. --- \VUgras{ખેડૂત} \textsuperscript{(36)} આ બધાં કામ ચલાવે છે અને
પોતાના નોકરોને સૂચના આપે છે. તે \VUgras{સમાર} \textsuperscript{(40)}
અને \VUgras{કોદાળી} \textsuperscript{(39)} છાપરા નીચે મૂકવાની ભલામણ કરે
છે, જે \VUgras{ઝૂંપડી} \textsuperscript{(38)} પાસે
ભુલાઈ ગયાં છે, જે \VUgras{કૂતરાની} \textsuperscript{(37)} છે.

%%K t07-c2-09-1 p
9. --- તેની બૂમોથી \VUgras{કબૂતર} \textsuperscript{(41)} ડરી જાય છે,
જે પૂરી પાંખે \VUgras{કબૂતરખાનામાં} \textsuperscript{(42)} ઊડી જાય છે,
અને \VUgras{મરઘીઓ} \textsuperscript{(43)} જે \VUgras{મરઘાંખાનામાં}
\textsuperscript{(44)} પાછી જવાની ઉતાવળ કરે છે.

%%K t07-c2-10-1 p
10. --- \VUgras{ઘેટાંના વાડા} \textsuperscript{(48)} આગળ, આ રહ્યો વૃદ્ધ
\VUgras{ભરવાડ} \textsuperscript{(45)} જે પોતાનું \VUgras{ટોળું}
\textsuperscript{(49)} પાછું લાવે છે. પોતાની \VUgras{લાકડી}
\textsuperscript{(46)} પકડીને, જે તેને કદી ખૂટતી નથી, પોતાના રંગ
ઊડેલા \VUgras{ઝભલાની} \textsuperscript{(47)} જેમ, તે વાડામાં
\VUgras{ઘેટાં} \textsuperscript{(50)}, \VUgras{ઘેટીઓ}
\textsuperscript{(53)}, \VUgras{ઘેટાંનાં બચ્ચાં} \textsuperscript{(54)},
\VUgras{બકરીઓ} \textsuperscript{(51)} અને \VUgras{બકરીનાં બચ્ચાં}
\textsuperscript{(52)} લઈ જાય છે. તેનો વફાદાર \VUgras{કૂતરો}
\textsuperscript{(55)} આર્ગુસ છેલ્લે આવે છે અને પોતાના ભસવાથી
મોડાં પડનારાંને ઉતાવળ કરાવે છે.

%%K t07-c2-11-1 p
11. --- આ બધો ઘોંઘાટ અને હલચલ સારી \VUgras{દૂધાળ ગાયની}
\textsuperscript{(56)} શાંતિ ડહોળતાં નથી, જે પોતાની \VUgras{ઘંટડી}
\textsuperscript{(57)} રણકાવે છે, જ્યારે \VUgras{ગમાણના}
\textsuperscript{(58)} બારણા આગળ તેને દોહવામાં આવે છે.

%%K t07-c2-12-1 p
12. તે સમજતી હોય એમ લાગે છે કે તેણે પોતાનું દૂધ આપવું જોઈએ, જેથી
\VUgras{ખેડૂતાણી} \textsuperscript{(60)} \VUgras{વલોણામાં}
\textsuperscript{(61)} \VUgras{માખણ} અથવા ઉત્તમ \VUgras{પનીર}
\textsuperscript{(64)} તૈયાર કરી શકે, જે \VUgras{મોટી કૂંડી}
\textsuperscript{(63)} પર મૂકેલા પાટિયામાંથી ટપકે છે.

%%K t07-c2-13-1 p
13. --- આખું \VUgras{દૂધઘર} \textsuperscript{(59)} \VUgras{ખેતરના
નોકર} \textsuperscript{(65)} દ્વારા ઘણું સ્વચ્છ રખાય છે, જેણે હમણાં જ
મોટી ડોલમાં \VUgras{દૂધનાં વાસણ} \textsuperscript{(62)} ધોયાં, જે તે
હાથમાં ઊંચકે છે.

%%K t07-tit-9 sub
\VUsaut{3.0mm}
\VUpk{10.2pt}{40}{મરઘાંખાનું.}
\VUsaut{3.0mm}

%%K t07-c2-14-1 p
14. --- પોતાનાં લાકડાનાં જાડાં ચંપલથી તે થોડું ભારે ચાલે છે, અને એમ
\VUgras{ટર્કી} \textsuperscript{(68)}, \VUgras{બતકીઓ}
\textsuperscript{(66)}, \VUgras{બતકો} \textsuperscript{(67)} અને
\VUgras{હંસ} \textsuperscript{(69)} ને ભગાડે છે, જે પોતાની
\VUgras{ચાંચ} \textsuperscript{(70)} પહોળી ખોલે છે અને બેચેન થઈને
ચીસો પાડે છે.

%%K t07-c2-14-2 p
\VUgras{કૂકડો} \textsuperscript{(71)}, વધારે લડાયક, પોતાની લાલ
\VUgras{કલગી} \textsuperscript{(72)} ટટ્ટાર કરે છે, પોતાની
\VUgras{પૂંછડીનાં} \textsuperscript{(73)} પીંછાં ખંખેરે છે અને
« કૂકડેકૂક » મોટેથી બોલે છે.

%%K t07-c2-15-1 p
15. --- \VUgras{માતા મરઘી} \textsuperscript{(74)} \VUgras{ખાતર}
\textsuperscript{(81)} પર ફરે છે પોતાનાં \VUgras{બચ્ચાં}
\textsuperscript{(75)} સાથે. \VUgras{ભૂંડનું બચ્ચું}
\textsuperscript{(78)} પોતાની મા તરફ ભાગે છે, પ્રચંડ \VUgras{ભૂંડણ}
\textsuperscript{(77)} જેને આપણે ભૂંડખાના પાસે જોઈએ છીએ.

%%K t07-c2-16-1 p
16. --- પોતાનાં ખાનાંમાં, \VUgras{સસલાં} \textsuperscript{(79)}
લોભથી નાજુક \VUgras{ઘાસ} \textsuperscript{(80)} ખાય છે જે ખેડૂતની દીકરી
તેમને લાવે છે. યોહાનિના પોતાનાં સસલાંને ઘણો પ્રેમ કરે છે અને, રોજ,
મરઘાંખાનામાં તાજાં \VUgras{ઈંડાં} \textsuperscript{(76)} લીધા પછી, તે
પોતાનાં નાનાં મિત્રોને મળવા આવે છે.

\VUsaut{6.0mm}
\VUfilet{22mm}
